\documentclass[12pt,a4paper]{report}
\usepackage[french]{babel}
\usepackage[utf8]{inputenc}
\usepackage[T1]{fontenc}
\usepackage{graphicx}
\usepackage{geometry}
\usepackage{xcolor}
\usepackage{tikz}
\usepackage{enumitem}
\usepackage{fancyhdr}
\setlength{\headheight}{14.5pt}
\addtolength{\topmargin}{-2.5pt}
\usepackage{hyperref}
\usepackage{tabularx}
\usepackage{array}
\usepackage{booktabs}
\setlength{\tabcolsep}{4pt}
\usepackage{listings}
\usepackage{mdframed}

\geometry{margin=2.5cm}

\definecolor{primary}{RGB}{41, 128, 185}
\definecolor{secondary}{RGB}{52, 73, 94}
\definecolor{accent}{RGB}{231, 76, 60}
\definecolor{lightgray}{RGB}{236, 240, 241}

\hypersetup{
    colorlinks=true,
    linkcolor=primary,
    urlcolor=primary
}

\pagestyle{fancy}
\fancyhf{}
\fancyhead[L]{\textbf{PEKEGNO} -- Plan de Découpage}
\fancyhead[R]{\thepage}

\lstset{
    basicstyle=\ttfamily\small,
    backgroundcolor=\color{lightgray},
    frame=single,
    breaklines=true
}

\begin{document}

\begin{titlepage}
    \centering
    \vspace*{3cm}
    {\Huge \textbf{PEKEGNO}}\\[1cm]
    {\Large Plateforme Web de Gestion Intégrée des Activités}\\[1.5cm]
    {\color{primary}\hrule height 3pt}
    \vspace{1cm}
    {\LARGE \textbf{Plan de Découpage du Projet}}\\[0.5cm]
    {\Large Organisation, Technologies et Répartition des Tâches}\\[2cm]
    \vfill
    \begin{tabularx}{0.6\textwidth}{>{\raggedleft\arraybackslash}X >{\raggedright\arraybackslash}X}
        \textbf{Durée} & 3 mois (90 jours) maximum \\
        \textbf{Version} & 1.0 \\
        \textbf{Date} & juillet 2026 \\
    \end{tabularx}
    \vspace{2cm}
\end{titlepage}

\tableofcontents
\newpage

% ============================================================================
\chapter{Présentation du Projet}
% ============================================================================

\section{Contexte}
\textbf{PEKEGNO} est une entreprise multi-agence proposant services, formations et activités commerciales. L'objectif est de développer une plateforme web SaaS centralisée pour gérer l'intégralité des opérations : CRM, formations (présentielles et distancielles), commercial, comptabilité, facturation, recouvrement et pilotage décisionnel.

\section{Durée et Organisation}
\begin{itemize}
    \item \textbf{Durée totale :} 3 mois calendaires maximum (90 jours)
    \item \textbf{Méthodologie :} Agile (Scrum) -- itérations de 5 à 7 jours
    \item \textbf{Validation :} Validation module par module avec le commanditaire
\end{itemize}

\section{Profils Utilisateurs}
\begin{enumerate}
    \item Super Administrateur
    \item Direction Générale
    \item Responsable Agence
    \item Responsable Département
    \item Commercial
    \item Caissier
    \item Comptable
    \item Formateur
    \item Client / Apprenant
\end{enumerate}

% ============================================================================
\chapter{Choix Technologiques}
% ============================================================================

\section{Architecture Globale}
\begin{center}
\begin{tikzpicture}[node distance=2cm, auto]
    \tikzstyle{block} = [rectangle, draw, fill=blue!10, rounded corners, text centered, minimum width=3cm, minimum height=1cm];
    \tikzstyle{arrow} = [->, thick];

    \node (client) [block] {Client (Web)};
    \node (api) [block, below of=client, yshift=-1cm] {API REST};
    \node (backend) [block, below of=api, yshift=-1cm] {Backend Laravel};
    \node (db) [block, below of=backend, yshift=-1cm] {Base de Données PostgreSQL};

    \draw[arrow] (client) -- node[right] {HTTPS} (api);
    \draw[arrow] (api) -- (backend);
    \draw[arrow] (backend) -- (db);
\end{tikzpicture}
\end{center}

\section{Stack Technique}

{\small
\noindent\begin{tabularx}{\textwidth}{l X}
    \toprule
    \textbf{Couche} & \textbf{Technologie} \\
    \midrule
    Frontend & React.js (Vite) + TypeScript + Tailwind CSS \\
    Backend & Laravel 11 (PHP 8.2+) \\
    Base de données & PostgreSQL 16 \\
    API & REST (JSON) + documentation Swagger/OpenAPI \\
    Authentification & JWT (Laravel Sanctum) + OAuth2 \\
    Paiements & MTN Mobile Money, Orange Money, Carte bancaire, PayPal \\
    Streaming vidéo & HLS sécurisé (AWS S3 + CloudFront) \\
    Stockage fichiers & AWS S3 (ou équivalent) \\
    Notifications & Email (Laravel Mail), SMS (option), WhatsApp (option) \\
    Hébergement & AWS / OVH / Scaleway \\
    CI/CD & GitHub Actions \\
    \bottomrule
\end{tabularx}}

\section{Base de Données}
\begin{itemize}
    \item \textbf{SGBD :} PostgreSQL 16
    \item \textbf{Justification :} Robustesse, support JSON, requêtes complexes, intégrité référentielle, excellent pour ERP/CRM multi-agence
    \item \textbf{Outil de migration :} Laravel Migrations
    \item \textbf{Approche :} Modélisation entité-relation normalisée (3NF) avec schémas par module
\end{itemize}

% ============================================================================
\chapter{Découpage du Projet -- Modules}
% ============================================================================

Le projet est découpé en \textbf{19 modules fonctionnels} répartis en \textbf{6 phases} de développement.

\section{Phase 1 -- Socle Technique (Jours 1--15)}
\label{sec:phase1}

\subsection*{Module 1 : Gestion des Agences}
\textbf{Durée :} 5 jours \\
\textbf{Description :} CRUD agences, informations (nom, pays, adresse, téléphone, email, responsable, statut), archivage.

\subsection*{Module 2 : Gestion des Départements}
\textbf{Durée :} 3 jours \\
\textbf{Description :} CRUD départements rattachés à une agence, affectation des utilisateurs et services.

\subsection*{Module 3 : Gestion des Catégories}
\textbf{Durée :} 3 jours \\
\textbf{Description :} CRUD catégories (nom, description, couleur, icône).

\subsection*{Module 7 : Gestion des Utilisateurs}
\textbf{Durée :} 5 jours \\
\textbf{Description :} CRUD utilisateurs (nom, pseudo, mot de passe, rôle, agence, département).

\subsection*{Module 8 : Gestion des Rôles}
\textbf{Durée :} 3 jours \\
\textbf{Description :} CRUD rôles (Direction, Admin, Responsable, Commercial, Comptable, Formateur, Caissier).

\subsection*{Module 9 : Gestion des Privilèges}
\textbf{Durée :} 3 jours \\
\textbf{Description :} Permissions par rôle : créer, modifier, supprimer, consulter, exporter, imprimer, valider, encaisser, annuler.

\section{Phase 2 -- Cœur Métier (Jours 16--48)}
\label{sec:phase2}

\subsection*{Module 4 : Gestion des Services / Activités}
\textbf{Durée :} 5 jours \\
\textbf{Description :} CRUD services (nom, catégorie, couverture, prix, description, vidéo, département, agence), promotions temporaires, historique des prix.

\subsection*{Module 5 : Gestion des Formations}
\textbf{Durée :} 12 jours \\
\textbf{Description :} Gestion des formations (présentiel/distanciel), modules pédagogiques, classement, sécurisation des contenus (streaming HLS, PDF protégé), paiements formation, suivi financier des apprenants.

\subsection*{Module 6 : Gestion des Formateurs}
\textbf{Durée :} 3 jours \\
\textbf{Description :} CRUD formateurs, planning, emploi du temps, affectation aux formations.

\section{Phase 3 -- CRM \& Commercial (Jours 49--63)}
\label{sec:phase3}

\subsection*{Module 10 : Gestion des Commerciaux}
\textbf{Durée :} 5 jours \\
\textbf{Description :} CRUD commerciaux, calcul automatique des commissions (\%/montant fixe/paliers), gestion des points (score, classement, dégressivité), historique commercial.

\subsection*{Module 11 : Gestion CRM}
\textbf{Durée :} 5 jours \\
\textbf{Description :} Gestion des prospects (statuts : nouveau, relance, devis envoyé, gagné, perdu), historique des interactions.

\subsection*{Module 12 : Vente / Commandes}
\textbf{Durée :} 5 jours \\
\textbf{Description :} Processus de vente, génération facture/créance/statistiques, multi-moyens de paiement.

\subsection*{Module 13 : Rapport Journalier}
\textbf{Durée :} 3 jours \\
\textbf{Description :} Rapport commercial quotidien, pénalité automatique en cas de non-saisie.

\section{Phase 4 -- Clients \& Facturation (Jours 64--75)}
\label{sec:phase4}

\subsection*{Module 14 : Gestion Clients}
\textbf{Durée :} 5 jours \\
\textbf{Description :} CRUD clients, espace formation (login/mot de passe), historique formations, paiements, créances, relances.

\subsection*{Module 15 : Facturation}
\textbf{Durée :} 5 jours \\
\textbf{Description :} Devis, factures, reçus (logo, entête, cachet, numérotation auto, QR code), statuts (brouillon, validée, partiellement payée, payée, annulée).

\subsection*{Module 16 : Gestion des Créances}
\textbf{Durée :} 5 jours \\
\textbf{Description :} Liste des créances, suivi (montant, échéance, retard, historique), relances automatiques.

\section{Phase 5 -- Comptabilité \& Reporting (Jours 76--84)}
\label{sec:phase5}

\subsection*{Module 17 : Comptabilité}
\textbf{Durée :} 8 jours \\
\textbf{Description :} Recettes, dépenses, catégories, balance, journal de caisse, trésorerie, états financiers.

\subsection*{Module 18 : Notifications}
\textbf{Durée :} 3 jours \\
\textbf{Description :} Notifications automatiques (Email, SMS, WhatsApp) \\ pour facture, échéance, relance, formation, paiement, rapport.

\subsection*{Module 19 : Tableau de Bord}
\textbf{Durée :} 8 jours \\
\textbf{Description :} Indicateurs (CA, trésorerie, créances, clients, apprenants, taux conversion), graphiques, alertes. Vues par profil (Direction, Commercial, Caissier, Formateur).

\section{Phase 6 -- Finalisation (Jours 85--90)}
\label{sec:phase6}
\begin{itemize}
    \item Tests d'intégration et recette générale
    \item Correction des bugs
    \item Déploiement production
    \item Documentation utilisateur
    \item Formation des utilisateurs clés
\end{itemize}

% ============================================================================
\chapter{Backend}
% ============================================================================

\section{Architecture Backend (Laravel 11)}
\begin{itemize}
    \item \textbf{Pattern :} MVC + Repository Pattern
    \item \textbf{API :} RESTful avec Laravel Sanctum (JWT)
    \item \textbf{Middleware :} Permission middleware basé sur les rôles et privilèges
    \item \textbf{Jobs/Queues :} Laravel Horizon pour les tâches asynchrones (notifications, relances, rapports)
    \item \textbf{Scheduler :} Tâches planifiées (sauvegardes, dégressivité des points, pénalités)
\end{itemize}

\section{Base de Données}
\begin{itemize}
    \item PostgreSQL 16 avec schémas par module
    \item Migrations Laravel versionnées
    \item Seeders pour les données de base (rôles, privilèges, admin)
    \item Indexation optimisée pour les performances
\end{itemize}

\section{Sécurité}
\begin{itemize}
    \item HTTPS obligatoire
    \item Authentification sécurisée (JWT + OAuth2)
    \item Double authentification (optionnelle)
    \item Blocage connexion multi-appareils
    \item Déconnexion automatique après inactivité
    \item Journalisation des connexions et actions
    \item Conformité OWASP
\end{itemize}

% ============================================================================
\chapter{Frontend}
% ============================================================================

\section{Architecture Frontend (React.js)}
\begin{itemize}
    \item \textbf{Framework :} React 18+ avec TypeScript
    \item \textbf{Build tool :} Vite
    \item \textbf{UI Library :} Tailwind CSS + Headless UI / Radix UI
    \item \textbf{State management :} Zustand ou Redux Toolkit
    \item \textbf{Routing :} React Router v6
    \item \textbf{Formulaires :} React Hook Form + Zod
    \item \textbf{Graphiques :} Recharts / Chart.js
    \item \textbf{API client :} Axios + React Query (TanStack Query)
    \item \textbf{Responsive :} Mobile-first, supports tablette et smartphone
\end{itemize}

\section{Expérience Utilisateur}
\begin{itemize}
    \item \textbf{Design :} Interface moderne, épurée, cohérente
    \item \textbf{Temps de réponse :} Inférieur à 3 secondes
    \item \textbf{Feedback :} États de chargement (skeleton), notifications toast, animations fluides
    \item \textbf{Accessibilité :} Navigation clavier, contrastes, labels ARIA
    \item \textbf{Performances :} Lazy loading, code splitting, pagination, filtres et recherche multicritère
\end{itemize}

% ============================================================================
\chapter{Stratégie de Test}
% ============================================================================

\section{Approche Générale}
Chaque module est testé immédiatement après son développement, avant de passer au suivant.

\section{Niveaux de Test}
\begin{enumerate}
    \item \textbf{Tests Unitaires} -- PHPUnit (backend), Vitest (frontend) \\
    Validation des fonctions et méthodes critiques
    \item \textbf{Tests d'API} -- Laravel Tests / Pest + Postman \\
    Validation des endpoints REST
    \item \textbf{Tests d'Intégration} -- Cypress / Playwright \\
    Parcours utilisateur complets
    \item \textbf{Tests de Sécurité} -- Analyse OWASP, test des authentifications
    \item \textbf{Tests de Performance} -- Vérification du temps de réponse < 3s
    \item \textbf{Tests de Recette} -- Validation par le commanditaire après chaque module
\end{enumerate}

\section{Procédure par Module}
\begin{mdframed}[backgroundcolor=lightgray!30, linecolor=primary]
\textbf{Processus :} Développement $\rightarrow$ Tests unitaires $\rightarrow$ Tests API $\rightarrow$ Validation interne $\rightarrow$ Recette commanditaire $\rightarrow$ Passage au module suivant
\end{mdframed}

% ============================================================================
\chapter{Répartition des Tâches}
% ============================================================================

\section{Équipe de Développement}

{\small
\noindent\begin{tabularx}{\textwidth}{l X}
    \toprule
    \textbf{Rôle} & \textbf{Responsabilités} \\
    \midrule
    Chef de Projet / Lead Dev & Architecture, coordination, validation technique, gestion des déploiements \\
    Développeur Backend (x1) & Laravel, API REST, base de données, intégrations paiement \\
    Développeur Frontend (x1) & React.js, UI/UX, intégration API, composants réutilisables \\
    Développeur Fullstack (x1) & Modules transverses, notifications, reporting, tableaux de bord \\
    \bottomrule
\end{tabularx}}

\section{Planning d'Exécution}

\begin{tabularx}{\textwidth}{l l X}
    \toprule
    \textbf{Phase} & \textbf{Jours} & \textbf{Tâches} \\
    \midrule
    Phase 1 & 1--15 & Socle technique : agences, départements, catégories, utilisateurs, rôles, privilèges \\
    Phase 2 & 16--48 & Services, formations, formateurs \\
    Phase 3 & 49--63 & CRM, commerciaux, ventes, rapports journaliers \\
    Phase 4 & 64--75 & Clients, facturation, créances \\
    Phase 5 & 76--84 & Comptabilité, notifications, tableaux de bord \\
    Phase 6 & 85--90 & Finalisation, tests, déploiement \\
    \bottomrule
\end{tabularx}

\section{Répartition Détaillée}

\subsection*{Développeur Backend}
\begin{itemize}
    \item Modélisation et implémentation de la base de données (PostgreSQL)
    \item API REST Laravel (tous les modules)
    \item Authentification JWT + gestion des sessions
    \item Intégration des moyens de paiement
    \item Système de notifications (Email, SMS, WhatsApp)
    \item Jobs asynchrones et tâches planifiées
    \item Documentation API (Swagger)
\end{itemize}

\subsection*{Développeur Frontend}
\begin{itemize}
    \item Maquettage et design de l'interface
    \item Composants React réutilisables
    \item Intégration des API (Axios + React Query)
    \item Pages responsive (ordinateur, tablette, smartphone)
    \item Graphiques et tableaux de bord
    \item Gestion des états (loading, empty, error)
\end{itemize}

\subsection*{Développeur Fullstack}
\begin{itemize}
    \item Modules 5 (Formations) et 12 (Ventes) -- les plus complexes
    \item Streaming sécurisé HLS
    \item Système de facturation et QR code
    \item Module de comptabilité
    \item Import/Export Excel/CSV
    \item Tests d'intégration (Cypress/Playwright)
\end{itemize}

\subsection*{Chef de Projet / Lead Dev}
\begin{itemize}
    \item Validation des spécifications
    \item Coordination entre backend et frontend
    \item Revue de code
    \item Gestion des déploiements (CI/CD)
    \item Communication avec le commanditaire
    \item Présentation des validations module par module
\end{itemize}

% ============================================================================
\chapter{Plan de Travail Progressif}
% ============================================================================

\section{Méthodologie de Validation}
Conformément aux instructions du commanditaire, le projet avance progressivement avec validation module par module :
\begin{enumerate}
    \item Développement du module
    \item Tests internes (unitaires + API)
    \item Présentation au commanditaire
    \item Validation ou feedback
    \item Passage au module suivant après validation
\end{enumerate}

\section{Ordre de Démarrage}
\textbf{Phase 1 -- Socle technique} (Jours 1--15) \\
Modules : Agences $\rightarrow$ Départements $\rightarrow$ Catégories $\rightarrow$ Utilisateurs $\rightarrow$ Rôles $\rightarrow$ Privilèges

Cet ordre permet de disposer immédiatement d'une base fonctionnelle (auth, permissions, structure multi-agence) sur laquelle tous les autres modules s'appuieront.

% ============================================================================
\chapter{Annexes}
% ============================================================================

\section{Glossaire}
\begin{tabularx}{\textwidth}{l X}
    \toprule
    \textbf{Terme} & \textbf{Définition} \\
    \midrule
    SaaS & Software as a Service \\
    CRM & Customer Relationship Management \\
    LMS & Learning Management System \\
    ERP & Enterprise Resource Planning \\
    JWT & JSON Web Token \\
    HLS & HTTP Live Streaming \\
    CRUD & Create, Read, Update, Delete \\
    \bottomrule
\end{tabularx}

\section{Estimation de l'Envergure}
PEKEGNO est un \textbf{ERP/CRM métier complet} combinant :
\begin{itemize}
    \item CRM (prospects et clients)
    \item LMS (formations en ligne)
    \item ERP commercial
    \item Gestion multi-agences
    \item Comptabilité et trésorerie
    \item Facturation et recouvrement
    \item Gestion documentaire
    \item Business Intelligence (tableaux de bord et reporting)
\end{itemize}

\end{document}
